% =====================================================================
% Smriti: Verifiable Agent Memory
% Paper draft skeleton — LaTeX
%
% Build:
%   pdflatex paper && bibtex paper && pdflatex paper && pdflatex paper
%
% Each section in this skeleton is a stub `\input{...}` that pulls from
% docs/papers/sections/. Drop the Markdown drafts (after pandoc
% conversion: `pandoc -f gfm -t latex section-N.md -o sections/N.tex`)
% into that directory.
%
% Bibliography: see refs.bib. Keys are the same as the [CITE: ...]
% placeholders in the Markdown drafts.
%
% Style: chosen to be portable. Uses generic `article` class with
% packages that produce conference-style output. Switch to `neurips_2026`,
% `usenix-2026`, or AAAI/IEEE styles for camera-ready by replacing
% \documentclass and minor preamble adjustments.
% =====================================================================

\documentclass[11pt, letterpaper]{article}

% --- core packages -------------------------------------------------------
\usepackage[margin=1in]{geometry}
\usepackage[T1]{fontenc}
\usepackage[utf8]{inputenc}
\usepackage{microtype}
\usepackage{lmodern}

% --- math ----------------------------------------------------------------
\usepackage{amsmath, amssymb, amsthm, mathtools}
\usepackage{bm}

% --- references and links ------------------------------------------------
\usepackage[
  backend=bibtex,
  style=numeric-comp,
  sorting=none,
  maxbibnames=99
]{biblatex}
\addbibresource{refs.bib}

\usepackage{hyperref}
\hypersetup{
  colorlinks=true,
  linkcolor=black,
  citecolor=blue!50!black,
  urlcolor=blue!50!black
}
\usepackage{cleveref}

% --- code listings -------------------------------------------------------
\usepackage{listings}
\usepackage{xcolor}
\definecolor{codebg}{RGB}{248,248,248}
\definecolor{codekey}{RGB}{10, 60, 120}
\definecolor{codecomment}{RGB}{90, 110, 90}
\lstset{
  basicstyle=\ttfamily\small,
  backgroundcolor=\color{codebg},
  keywordstyle=\color{codekey}\bfseries,
  commentstyle=\color{codecomment}\itshape,
  frame=single,
  framesep=4pt,
  framerule=0.4pt,
  rulecolor=\color{black!30},
  showstringspaces=false,
  breaklines=true,
  captionpos=b,
  numberstyle=\tiny\color{black!50},
  numbers=left,
  numbersep=8pt,
  xleftmargin=2.5em
}
\lstdefinelanguage{Rust}{
  morekeywords={fn, let, mut, pub, struct, impl, trait, async, await, return, use, mod, self, Self, match, if, else, for, in, while, loop, break, continue, where, type, enum, const, static, ref, dyn, as, move, unsafe, crate, super, true, false},
  morecomment=[l]{//},
  morecomment=[s]{/*}{*/},
  morestring=[b]"
}

% --- tables --------------------------------------------------------------
\usepackage{booktabs}
\usepackage{array}
\usepackage{multirow}

% --- algorithms ----------------------------------------------------------
\usepackage{algorithm}
\usepackage{algpseudocode}

% --- diagrams ------------------------------------------------------------
\usepackage{tikz}
\usetikzlibrary{shapes, arrows.meta, positioning, fit, backgrounds, calc}

% --- theorem environments ------------------------------------------------
\theoremstyle{definition}
\newtheorem{definition}{Definition}[section]
\newtheorem{invariant}[definition]{Invariant}

\theoremstyle{plain}
\newtheorem{theorem}[definition]{Theorem}
\newtheorem{lemma}[definition]{Lemma}

\theoremstyle{remark}
\newtheorem*{remark}{Remark}

% --- author shortcuts ----------------------------------------------------
\newcommand{\Smriti}{\textsc{Smriti}}
\newcommand{\sysname}{\Smriti}
\newcommand{\Iname}{\textbf{I1}}
\newcommand{\IIname}{\textbf{I2}}
\newcommand{\IIIname}{\textbf{I3}}
\newcommand{\sha}{\mathrm{SHA\text{-}256}}
\newcommand{\overlap}{\omega}

% --- title ---------------------------------------------------------------
\title{
  Verifiable Agent Memory: \\
  A Hash-Chained Integrity Layer for LLM-Generated Knowledge Graphs
}

\author{
  Anonymous Authors \\
  \texttt{anonymous@example.org}
}

\date{\today}

% =====================================================================
\begin{document}

\maketitle

\begin{abstract}
\input{sections/abstract}
\end{abstract}

% --- §1 -----------------------------------------------------------------
\section{Introduction}
\label{sec:intro}
\input{sections/01-introduction}

% --- §2 -----------------------------------------------------------------
\section{Background}
\label{sec:background}
% TODO: draft 02-background. Convert from docs/papers/sections/<corresponding-md>.md via pandoc, then drop here.


% --- §3 -----------------------------------------------------------------
\section{The Integrity Contract}
\label{sec:contract}
\input{sections/03-integrity-contract}

% --- §4 -----------------------------------------------------------------
\section{Architecture}
\label{sec:architecture}
% TODO: draft 04-architecture. Convert from docs/papers/sections/<corresponding-md>.md via pandoc, then drop here.


% --- §5 -----------------------------------------------------------------
\section{Implementation}
\label{sec:implementation}
% TODO: draft 05-implementation. Convert from docs/papers/sections/<corresponding-md>.md via pandoc, then drop here.


% --- §6 -----------------------------------------------------------------
\section{Evaluation}
\label{sec:evaluation}

\subsection{Audit overhead}
\label{sec:eval-overhead}
\input{sections/06-1-audit-overhead}

\subsection{Chain integrity at scale}
\label{sec:eval-chain}
% TODO: draft 06-2-chain-scale. Convert from docs/papers/sections/<corresponding-md>.md via pandoc, then drop here.


\subsection{Reproducibility-by-replay}
\label{sec:eval-replay}
% TODO: draft 06-3-replay-reproducibility. Convert from docs/papers/sections/<corresponding-md>.md via pandoc, then drop here.


\subsection{Provenance threshold sweep}
\label{sec:eval-factum}
% TODO: draft 06-4-factum-threshold. Convert from docs/papers/sections/<corresponding-md>.md via pandoc, then drop here.


\subsection{Worked replay case study}
\label{sec:eval-case-study}
% TODO: draft 06-5-replay-case-study. Convert from docs/papers/sections/<corresponding-md>.md via pandoc, then drop here.


% --- §7 -----------------------------------------------------------------
\section{Related Work}
\label{sec:related}
% TODO: draft 07-related-work. Convert from docs/papers/sections/<corresponding-md>.md via pandoc, then drop here.


% --- §8 -----------------------------------------------------------------
\section{Limitations and Future Work}
\label{sec:limitations}
% TODO: draft 08-limitations. Convert from docs/papers/sections/<corresponding-md>.md via pandoc, then drop here.


% --- §9 -----------------------------------------------------------------
\section{Conclusion}
\label{sec:conclusion}
% TODO: draft 09-conclusion. Convert from docs/papers/sections/<corresponding-md>.md via pandoc, then drop here.


% --- bibliography --------------------------------------------------------
\printbibliography

% --- appendices ----------------------------------------------------------
\appendix

\section{Schema reference}
\label{app:schema}
% TODO: draft appendix-A-schema. Convert from docs/papers/sections/<corresponding-md>.md via pandoc, then drop here.


\section{Reproducibility checklist}
\label{app:repro}
% TODO: draft appendix-B-reproducibility. Convert from docs/papers/sections/<corresponding-md>.md via pandoc, then drop here.


\section{Glossary}
\label{app:glossary}
% TODO: draft appendix-C-glossary. Convert from docs/papers/sections/<corresponding-md>.md via pandoc, then drop here.


\end{document}
